\section{评估方法}
\label{sec:method}

三个实验项目在同一个产品上先后进行（表~\ref{tab:programs}）。V1 在自建任务集上检验团队设计和效率；V8 通过 Harbor
在 Terminal-Bench~4.0 上检验各种手段；最后一次运行按 Artificial Analysis 跑智能体的方式，把选定的配置放到
Terminal-Bench~4.1 上。

\begin{table}[t]
\centering
\small
\caption{三个实验项目。V8 的试验数包括每一次已评分的尝试，含筛查和重跑。}
\label{tab:programs}
\begin{tabularx}{\textwidth}{@{}l l L >{\raggedright\arraybackslash}p{30mm} >{\raggedright\arraybackslash}p{33mm}@{}}
\toprule
项目 & 日期（2026） & 任务 & 模型 & 规模 \\
\midrule
V1 & 9/28 -- 10/26 & 自建任务集：自编任务、代码仓库任务、公开任务（Terminal-Bench、DeepSWE）和问答任务，都带隐藏
  评分器 & 先 GPT-6 Astra，后 GPT-6.1 Sol，中等强度 & 两个本地评估实例上 120 多批 \\
V8 & 10/27 -- 11/19 & Terminal-Bench~4.0：不需要 GPU 的 63 道题，分为可读 33 道、未读 30 道 & GPT-6.1 Sol、GPT-6.1
  Terra（筛查）、GPT-6.2 Sol & 一台 Linux 电脑上 315 个 Harbor 作业、6,223 次已评分试验 \\
最终 & 11/20 & Terminal-Bench~4.1，AA 的数据集摘要版本，63 道题 & GPT-6.2 Sol，max 强度 & 189 次尝试，每题 3 次 \\
\bottomrule
\end{tabularx}
\end{table}

\subsection{任务集}

\para{V1 任务集}每道题都带一个隐藏评分器和一份参考解。一道题进入任务集之前，先由一个全新的模型在空白工作区里盲做
一遍，用它的得分检验评分器；在早期的一批里，这样找出了两个真实的评分器缺陷。已经饱和或不公平的题会被撤下
（10 月 9 日撤下 52 道，剩 32 道在用）。一个设计在它所针对的题上胜出之后，还必须在不针对的题上站得住，才能上线。

\para{V8 的划分}10 月 30 日，63 道能运行的 Terminal-Bench~4.0 题目被分成 25 道开发题和 38 道留出题。第一次读取留出
结果（11 月 4 日）之后，榜单上有智能体胜过我们的 8 道留出题被打开研究。这样剩下 33 道\textbf{可读}题，我们阅读它们的
轨迹、失败的测试和题目文件来设计手段；另有 30 道\textbf{未读}题用于确认。未读题只以总和的形式出现：奖励、成本、
分钟数和模型调用次数的合计；它们的题名、轨迹和测试从未打开过，当时没有，写这篇论文时也没有。

\subsection{基础设施}

V8 使用 Harbor 0.24.0（最终运行用 0.26.0）\citep{harbor2026}，运行在一台 32 线程、96\,GB 内存的台式电脑上；在更高
负载下一次内存交换风暴损失了 44 次运行，此后同时运行约 30 次试验。团队自身的时间上限是 465 分钟，题目的上限是
480 分钟，因为服务端默认的 3 小时曾把一道需要 8 小时的题提前截断。11 月 5 日之前，运行在一台 Mac 上的 Docker
虚拟机里进行，每次同时跑两次试验；它的留出基线比那台电脑低，这也是每次比较都在同一批次、同一台机器上进行的原因
之一。

\subsection{指标}

题目 $t$ 的第 $d$ 次运行的奖励 $r_{t,d}$ 为 0 或 1。我们的得分是各题平均奖励之和：
\begin{equation}
  S = \sum_{t \in T} \frac{1}{n_t} \sum_{d=1}^{n_t} r_{t,d},
  \label{eq:score}
\end{equation}
因此 $S/|T|$ 就是通过率。缺失或没有奖励的尝试记为 0。成本是该次运行的 token 按服务端价格表中 OpenAI 的 API 公开标价
（2026 年 10 月 18 日核对）计算的，尽管多数运行实际使用的是 ChatGPT（Codex）账号登录。时间是智能体自己每次尝试所用的
时间；V1 的效率研究用的是到最后一次模型调用为止的活跃时间。

\subsection{统计分析}

在 Terminal-Bench 题集上，实验组和基线在每道题上各运行两次或更多。我们检验 $S$ 的差异时，在每道题内部把各次运行
在实验组和基线之间重新分配（20,000 次置换），由此得到零假设下的标准差和 $p$ 值。在未读的 30 道题上每边两次运行时，
零假设标准差约为 1.4（各次读取在 1.33 到 1.53 之间），因此一次读取能分辨约 2.7 道题的差异；每边四次运行时约为 0.94。
凡给出标准误的地方，标准误都是按题计算的。V1 对配对运行使用成本或时间对数比的精确 Wilcoxon 符号秩检验
\citep{wilcoxon1945}。

\subsection{实验规范}

\begin{itemize}[leftmargin=1.4em]
  \item \textbf{同批基线。}每个实验组都和它的基线在同一时段运行，因为曾有一个基线在两批之间移动了 +1.67。
  \item \textbf{空闲重评。}每个零分都在空闲机器上重评一次，因为负载下验证器超时会产生假零分。在空闲机器上这种情况
    很少：两次大规模重评中，43 个零分里 43 个、54 个零分里 54 个仍为零；第三次重评中，74 个里有 1 个变了。
  \item \textbf{预先登记。}每批实验之前，都把规则、分支和我们自己的预测写进实验记录本 \citep{nosek2018prereg}；V8 的
    记录本最终有 9,276 行、387 节。我们还请一个独立的审查智能体专门找设计和解读中的漏洞。
\end{itemize}
